\section{INTRODUCCIÓN}

La Web Semántica propone un cambio de enfoque en la publicación y recuperación de
información en la Web: los datos dejan de ser únicamente documentos legibles por
personas y pasan a representarse también como conocimiento formal interpretable por
máquinas \parencite{berners-lee2001}. Este enfoque se sustenta en estándares del W3C
como RDF para la representación de grafos de conocimiento, RDFS y OWL para la
definición de vocabularios y ontologías, y SPARQL para la consulta declarativa sobre
esos grafos.

En plataformas de alojamiento temporal tipo Airbnb, la necesidad de representar
conocimiento semántico es especialmente evidente. La decisión de reserva no depende de
un solo atributo, sino de combinaciones entre tipo de propiedad, rango de precio,
ubicación, amenidades, perfil del viajero y propósito del viaje. Una consulta como
``quiero un departamento tranquilo, con Wi-Fi, adecuado para trabajo remoto y cercano a
una universidad'' contiene relaciones implícitas entre conceptos que un buscador basado
solo en filtros discretos no modela ni explota adecuadamente.

En este contexto, el presente trabajo desarrolla un buscador semántico para el dominio
de alojamiento temporal. La solución se apoya en una ontología OWL/RDF propia cargada
en Apache Jena Fuseki, sobre la cual se ejecutan consultas SPARQL para recuperar y
ordenar resultados. El sistema no se limita a buscar coincidencias textuales, sino que
combina restricciones duras sobre precio, capacidad o ciudad con criterios semánticos
relacionados con amenidades, perfiles de huésped y propósito del viaje.

Un requisito central del enunciado fue la integración con DBpedia. En la solución
implementada, esta integración se materializa mediante el enriquecimiento de la
ontología local con enlaces \texttt{owl:sameAs} y \texttt{owl:equivalentClass} hacia
recursos de DBpedia vinculados con tipos de alojamiento, amenidades y ciudades de
Bolivia. Este enfoque permite incorporar conocimiento geográfico real y, al mismo
tiempo, mantener una demostración estable en entorno local sobre Fuseki sin depender de
la disponibilidad del \textit{endpoint} público en cada consulta.

Otro requisito explícito fue la multilingüidad. Por ello, la ontología fue enriquecida
con anotaciones \texttt{rdfs:label} y \texttt{rdfs:comment} en tres idiomas:
español, inglés y francés. Esta representación permite reutilizar un único grafo RDF
para distintas interfaces de usuario y demuestra un uso correcto de las capacidades de
internacionalización propias de RDF/OWL.

La aplicación resultante fue implementada con Next.js como capa de presentación y de
servicios, Fuseki como \textit{triple store}, Bun como entorno de ejecución y un módulo
de procesamiento lingüístico para interpretar consultas en lenguaje natural. El sistema final ofrece una
interfaz web con navegación multilingüe, resultados visuales por tarjetas y búsqueda
semántica sobre una ontología consistente y validada.

El informe se organiza en cinco bloques principales: introducción y objetivos del
trabajo, marco teórico de las tecnologías de Web Semántica utilizadas, marco
referencial del dominio seleccionado, sección de ingeniería con la descripción de la
ontología, la arquitectura y las pruebas del prototipo, y finalmente las conclusiones
obtenidas a partir del desarrollo realizado.
